\documentclass[10pt,a4paper]{amsart}
\usepackage[margin=19mm]{geometry}
\usepackage{amsmath,amssymb,mathtools}
\usepackage[scr=rsfs,cal=euler]{mathalfa}
\usepackage{xfrac}
\usepackage[unicode,hidelinks,bookmarks=false]{hyperref}
\newcommand{\ie}{i.e.\ }
\newcommand{\cf}{cf.\ }
\newcommand{\resp}{respectively\ }
\newcommand{\Subsection}{\subsection}
\providecommand{\nobreakdash}{\nobreak\hbox{-}}
\setlength{\emergencystretch}{2em}
\multlinegap0pt
\usepackage{bm}
\usepackage{bbm}
\usepackage{dsfont}
\usepackage{xspace}
\usepackage{cancel}
\usepackage{mathtools}
\usepackage{amsthm}
\makeatletter
\@ifundefined{lemma}{\newtheorem{lemma}{Lemma}}{}
\@ifundefined{theorem}{\newtheorem{theorem}{Theorem}}{}
\makeatother
\title[Similarity-based Outlier Detection for Noisy Object Re-Ident]{Similarity-based Outlier Detection for Noisy Object Re-Identification Using Beta Mixtures}
\author{Waqar Ahmad, Evan Murphy, Vladimir A. Krylov}
\date{Source: arXiv:2509.08926v3 (CC BY 4.0)}
\begin{document}
\maketitle

\noindent\emph{Source and licence.} Similarity-based Outlier Detection for Noisy Object Re-Identification Using Beta Mixtures. Version arXiv:2509.08926v3 (CC BY 4.0), distributed under the Creative Commons Attribution 4.0 International licence, as recorded in the arXiv licence metadata for this exact version and shown on the abstract page https://arxiv.org/abs/2509.08926v3. Full rights notice: licenses/arxiv-2509.08926-CC-BY-4.0.txt.\par\smallskip

\noindent\textit{Editorial scope.} Counted unit: the Theorem whose source label is \texttt{th:1} in the article (source0.tex lines 302--308, source label \texttt{th:1}), delivered together with the article's own complete original proof of it (source0.tex lines 310--369, one \texttt{proof} environment of 60 lines). No mathematics is written, repaired or re-derived: statement and proof are the article's own text line for line. Required context retained uncounted: lem:1 (lines 279--281). Every definition, display and label of those lines is retained unchanged. Omitted from this package and not redistributed: 1--278, 282--301, 309--309, 370--734. Nothing inside a retained excerpt is omitted or altered: every retained body is delivered exactly as the article's lines are written, so no omission receipt of any kind accompanies this package. Citations: the retained excerpts carry no citation command at all, so no bibliography entry of the article is transcribed and no reference is redistributed. Reference adaptation: none was needed and none was made; every reference inside the delivered unit resolves inside it, so no unresolved reference or citation is produced. Printed slips kept literal: no printed word, symbol, hypothesis or number of the counted statement or of its proof was altered. Layout and class: the article's own class is \texttt{IEEEtran} with packages including amsmath, amssymb, amsfonts, array, subfig, textcomp, stfloats, url, verbatim, graphicx, cite, booktabs, tabularray, algpseudocode; the wrapper is the lane's pinned \texttt{amsart} editorial wrapper at 10pt on A4, which restates the theorem-like environments the retained text needs and adds the article's own macro definitions that the retained text needs (none), with extra wrapper packages (bm, bbm, dsfont, xspace, cancel, mathtools). The delivered numbering is the unit's own. Assets: no retained span contains a figure environment, a graphics-inclusion command, a PDF-page inclusion, a table or a TikZ picture; no image file is copied into this package, the package declares no asset dependency and no image file is redistributed with it. Licence: version arXiv:2509.08926v3 of this article is distributed under the Creative Commons Attribution 4.0 International licence, as recorded in the arXiv licence metadata for this exact version and shown on the abstract page https://arxiv.org/abs/2509.08926v3. A full rights notice is delivered at licenses/arxiv-2509.08926-CC-BY-4.0.txt. Characters: every retained excerpt is pure ASCII, so the delivered engine, XeLaTeX with its default Latin Modern text font, reports no missing character.
\begin{lemma}\label{lem:1}
    For any $n\in\mathbb{Z}_{\geq 0}$, $\lim\limits_{\varepsilon\to 0^+} \varepsilon \Gamma(-n+\varepsilon) = \frac{(-1)^n}{n!}$.
\end{lemma}

\begin{theorem} \label{th:1}
    Let $(a_1,b_1),(a_2,b_2) \in \mathbb{R}^2_{>0}$, $\pi \in (0,1)$, and consider the mixture distribution
    \begin{align*}
        F(x) = \pi f(x;a_1,b_1)+(1-\pi)f(x;a_2,b_2).
    \end{align*}
    If $(a_1,b_1)-(a_2,b_2) \notin \mathbb{Z}^2$, or, in other words, either $a_1-a_2$ or $b_1-b_2$ (or both) are not integers, then $F$ has a unique representation as a mixture of two Beta distributions.
\end{theorem}

\begin{proof}
Let $\pi \in (0,1)$, and $(a_1,b_1),(a_2,b_2)\in \mathbb{R}^2_{>0}$ be distinct, such that $(a_1,b_1)-(a_2,b_2) \notin \mathbb{Z}^2$. Without loss of generality, assume that $a_1-a_2 \notin \mathbb{Z}$, and $a_1<a_2$. Indeed, if only $b_1-b_2 \notin \mathbb{Z}$ holds, we can then swap variable to $y:=1-x$ and revert to the $a_1-a_2 \notin \mathbb{Z}$ case due to the symmetry of the Beta distribution.

We will now use the integral Mellin transform defined for any $(\alpha,\beta)\in \mathbb{R}^2_{>0}$ as
\begin{align}
    m_{\alpha,\beta} = \int\limits_0^1 x^{t-1} f(x;\alpha,\beta)\,\text{d}x = \frac{\Gamma(t+\alpha-1)\Gamma(\alpha+\beta)}{\Gamma(t+\alpha+\beta-1)\Gamma(\alpha)}, \label{eq:mellin}
\end{align}
which can also be applied to the mixture $F$, yielding
\begin{align*}
    M_F(t) = \pi m_{a_1,b_1}(t)+(1-\pi)m_{a_2,b_2}.
\end{align*}
Note that the right hand side in Eq.~\ref{eq:mellin} has singularities for $t \in S_{\alpha,\beta}= \{1-\alpha-k:k \in \mathbb{Z}_{\geq 0}\}$ if $\beta \notin \mathbb{Z}_{>0}$, and $t \in S_{\alpha,\beta}= \{1-\alpha-k:k=0,...,\beta-1\}$ if $\beta \in \mathbb{Z}_{>0}$. To see this, notice that when $\beta \notin \mathbb{Z}_{>0}$, $\Gamma(t+\alpha-1)$ is non-finite if and only if $t \in S_{\alpha,\beta}$, and is the only non-finite term appearing in $S_{\alpha,\beta}$ in this case. Since the denominator is never zero due to the properties of $\Gamma$, these are the only values at which the numerator is non-finite, and a singularity exists. When $\beta \in \mathbb{Z}_{>0}$, using the core fact that $\Gamma(z+1)=z\Gamma(z)$ for all $z \in \mathbb{R}$, we can write 
\begin{align*}
\frac{\Gamma(1+\alpha-k)}{\Gamma(1+\alpha+\beta-k)} = \frac{1}{(1+\alpha-k)\cdots(1+\alpha+(\beta-1)-k)}.
\end{align*}
This expression has a singularity precisely when the denominator is zero, which happens if and only if $t \in S_{\alpha,\beta}$.

We now have that $m_{\alpha,\beta}$ converges outside of the set of singularities. Let $S:=S_{a_1,b_1} \cup S_{a_2,b_2}$ be the singularities of $M_F(t)$. Hence
$$
    \max S = \max S_{a_1,b_1} \cup S_{a_2,b_2}
    = \max \{1-a_1,1-a_2\} = 1-a_1.
$$
Thus, parameter $a_1$ can be uniquely recovered from observing the behavior of the mixture $F$ through $M_F(t)$. Moreover, as $S_{a_1,b_1}$ and $S_{a_2,b_2}$ are disjoint by the fact that $a_1-a_2 \notin \mathbb{Z}$, 
\begin{align*}
    \max S\backslash S_{a_1,b_1} = \max S_{a_2,b_2} = 1-a_2.
\end{align*}
Hence the parameter $a_2$ can also be recovered.

We can now address the recovery of $b_1,b_2$. Note that if $b_1 \in\mathbb{Z}_{>0}$ or $b_2 \in \mathbb{Z}_{>0}$, then $|S_{a_1,b_1}| = b_1$ or $|S_{a_2,b_2}| = b_2$, respectively. Also, we may recover $\pi$ by noticing that
\begin{align*}
    \pi = \frac{F(x)-f(x;a_2,b_2)}{f(x;a_1,b_1)-f(x;a_2,b_2)},
\end{align*}
and evaluating the right hand side at any $x \in (0,1)$. Going forward, we assume $(b_1,b_2) \notin \mathbb{Z}^2$. Define the sequence $\{t_n\}_{n\in\mathbb{Z}_{>0}}$, where $t_n = 1-a_1-n$. We now define residuals
$$R_n :=\lim\limits_{t \to t_n^+} (t-t_n)M_F(t).$$
Since $F$ is known (or observed), so is $M_F(t)$ for any $t$, and hence $R_n$ can be computed for each $n$. Moreover, as $m_{a_2,b_2}(t_n)$ is finite for each $n$, we have
\begin{align*}
    R_n &= \lim\limits_{t \to t_n^+} (t-t_n)[\pi m_{a_1,b_1}(t)+(1-\pi)m_{a_2,b_2}]\\
    &= \pi\lim\limits_{t \to t_n^+} (t-t_n)m_{a_1,b_1}(t).
\end{align*}
Making the substitution $\varepsilon = t_n-t$ and using Eq.~\ref{eq:mellin} we get
\begin{align*}
    &R_n = \pi \lim\limits_{\varepsilon \to 0^+} \varepsilon\frac{\Gamma(-n+\varepsilon)\Gamma(a_1+b_1)}{\Gamma(-n+b+\varepsilon)\Gamma(a_1)}\\
    &= \frac{\pi \Gamma(a_1+b_1)}{\Gamma(a_1)}\lim\limits_{\varepsilon \to 0^+} \left[\frac{1}{\Gamma(-n+b_1+\varepsilon)}\right] \lim\limits_{\varepsilon\to 0^+} \left[\varepsilon \Gamma(-n+\varepsilon)\right]  \\
    &=\frac{\pi \Gamma(a_1+b_1)}{\Gamma(a_1)} \frac{1}{\Gamma(-n+b_1)} \frac{(-1)^n}{n!},
\end{align*}
where the last equality uses Lemma \ref{lem:1}. For $n=0,1$, we have
\begin{align*}
    &R_0 = \frac{\pi\Gamma(a_1+b_1)}{\Gamma(a_1)\Gamma(b_1)},\\
    &R_1 = -\frac{\pi \Gamma(a_1+b_1)}{\Gamma(a_1)\Gamma(b_1-1)} = -\frac{\pi(b_1-1)\Gamma(a_1+b_1)}{\Gamma(a_1)\Gamma(b_1)},
\end{align*}
where the last equality follows from the fact that $\Gamma(b_1)=(b_1-1)\Gamma(b_1-1)$. Then notice that
\begin{align*}
    \frac{R_1}{R_0}  = -(b_1-1).
\end{align*}
Thus, we can recover $b_1$ and $\pi$ as
\begin{align*}
    b_1 = 1-\frac{R_1}{R_0}, \quad \pi = R_0 \frac{\Gamma(a_1)\Gamma(b_1)}{\Gamma(a_1+b_1)}.
\end{align*}
The recovery of $b_2$ follows an identical logic, using the sequence $t_n'=1-a_2-n$. Therefore, we have recovered unique values of $a_1,b_1,a_2,b_2$ and $\pi$. Thus, $F$ must have a unique representation as a mixture of two Beta distributions.
\end{proof}

\end{document}
